\clearpage
\thispagestyle{plain}
\oldpage[IV-4]{5}
\phantomsection\label{chapter:IV4-fr}
\begin{center}
{\large CAPÍTULO IV \textit{(fin)}\par}
\vspace{1.5em}
{\Large\bfseries ESTUDIO LOCAL DE LOS ESQUEMAS\\
Y DE LOS MORFISMOS DE ESQUEMAS\par}
\end{center}
\section{Invariantes diferenciales. Morfismos diferencialmente suaves}
\label{section:IV.16}

En este parágrafo presentamos, en forma global, algunas nociones de cálculo
diferencial particularmente útiles en Geometría algebraica. Omitimos numerosos
desarrollos clásicos de la Geometría diferencial (conexiones, transformaciones
infinitesimales asociadas a un campo de vectores, jets, etc.), aunque estas nociones
se expresan de manera particularmente natural en el marco de los esquemas. Omitimos
asimismo aquí los fenómenos propios de la característica $p>0$ (algunos de los cuales
se estudian, en el marco afín, en (\textbf 0, 21)). Para ciertos complementos sobre el
formalismo diferencial en los preesquemas, el lector podrá consultar las exposiciones
II y VII de [42], así como los capítulos posteriores de este Tratado.
